\chapter{Crédit et levier dans la finance décentralisée}\label{ch:m3:defi-credit-and-leverage}

Le 11 octobre 2022, un trader a acheté un \omterm{def:m3:blockchains-for-traders:key}{jeton} peu liquide sur les quelques bourses dont les cours alimentaient l'\omterm{def:m3:blockchains-for-traders:contract}{oracle} d'une plateforme de prêt.
En moins d'une demi-heure, le prix remonté par l'\omterm{def:m3:blockchains-for-traders:contract}{oracle} a été multiplié par plus de treize. La position du trader sur
la plateforme, désormais valorisée à ce prix, valait assez pour servir de collatéral à l'emprunt de tout ce que la plateforme
détenait~: il a emprunté et retiré plus de 110 millions de dollars d'autres \omterm{def:m3:blockchains-for-traders:key}{jetons} et a laissé
derrière lui le collatéral gonflé. Le prêt sur une \omterm{def:m3:blockchains-for-traders:chain}{blockchain} se passe de chargé de crédit~: un programme prête contre collatéral
selon un \omterm{def:m3:defi-credit-and-leverage:ltv}{ratio prêt/valeur} qu'il connaît, liquide quiconque dont le collatéral passe sous un seuil, et valorise tout
avec un \omterm{def:m3:blockchains-for-traders:contract}{oracle}. Ce chapitre traite des pools qui assurent les prêts, des liquidations qui les maintiennent solvables, des
\omterm{def:m3:defi-credit-and-leverage:flash}{flash loans} qui n'existent nulle part ailleurs, du \omterm{def:m3:defi-credit-and-leverage:staking}{rendement de staking} sur lequel repose une part importante du collatéral, et des
deux façons dont le système s'est brisé~: les \omterm{def:m3:blockchains-for-traders:stable}{stablecoins} dont l'arrimage reposait sur un algorithme, et les \omterm{def:m3:blockchains-for-traders:contract}{oracles}
qu'on pouvait acheter.

\section{Pools de prêt}

\begin{definition}[Pool de prêt, courbe de taux coudée]\label{def:m3:defi-credit-and-leverage:pool}
Un \emph{pool de prêt}\index{pool de prêt} est un \omterm{def:m3:blockchains-for-traders:contract}{smart contract} qui reçoit des dépôts d'un \omterm{def:m3:blockchains-for-traders:key}{jeton} de la part de fournisseurs et
les prête à des emprunteurs qui déposent d'autres \omterm{def:m3:blockchains-for-traders:key}{jetons} en collatéral, reversant aux fournisseurs les intérêts payés par les emprunteurs, diminués
d'une part conservée par le protocole. Une \emph{courbe de taux coudée}\index{courbe de taux coudée} fixe le taux d'emprunt
en fonction du taux d'utilisation du pool, croissant doucement jusqu'à un taux d'utilisation cible et fortement au-delà.
\end{definition}

Le taux d'utilisation est la part des \omterm{def:m3:blockchains-for-traders:key}{jetons} fournis qui est prêtée, comme dans le prêt de titres du livre~1,
chapitre~16. Le pool n'a pas d'échéance~: les fournisseurs peuvent retirer à tout moment, mais seulement ce qui n'est pas prêté. Le
coude est la défense du pool pour préserver cette liquidité. Au-delà, l'emprunt devient si cher que les emprunteurs remboursent et
les fournisseurs affluent, ramenant le taux d'utilisation à la baisse et laissant des \omterm{def:m3:blockchains-for-traders:key}{jetons} disponibles pour ceux qui veulent retirer.

\begin{proposition}[Taux servi aux fournisseurs]\label{prop:m3:defi-credit-and-leverage:supply}
Avec un taux d'utilisation $u$, un taux d'emprunt $r_b(u)$ et un facteur de réserve $f_{\mathrm{res}}$ conservé par le protocole, les fournisseurs perçoivent
$r_s = r_b(u)\,u\,(1 - f_{\mathrm{res}})$.
\end{proposition}

\begin{proof}
Les emprunteurs paient $r_b$ sur le montant prêté, soit une fraction $u$ des fonds fournis~; le protocole conserve $f_{\mathrm{res}}$ de ce montant.
\end{proof}

\begin{omfigure}
\begin{tikzpicture}
\begin{axis}[width=11cm,height=5cm,
  xlabel={taux d'utilisation, \%},ylabel={taux, \% par an},
  tick label style={font=\small},label style={font=\small},
  xmin=0,xmax=100,ymin=0,ymax=65,grid=major,grid style={gray!25},
  legend columns=2,legend style={font=\footnotesize,at={(0.5,-0.3)},anchor=north,draw=none,fill=none,/tikz/every even column/.append style={column sep=8pt}},
  table/col sep=comma]
\addplot[omThm,very thick] table[x=u,y=borrow]{figdata/markets-3/23-defi-credit-and-leverage/rates.csv};\addlegendentry{taux d'emprunt}
\addplot[omDef,very thick,dashed] table[x=u,y=supply]{figdata/markets-3/23-defi-credit-and-leverage/rates.csv};\addlegendentry{taux servi aux fournisseurs}
\end{axis}
\end{tikzpicture}

\omcaption{Une \omterm{def:m3:defi-credit-and-leverage:pool}{courbe de taux coudée}~: 4~\% au taux d'utilisation optimal de 90~\%, puis une hausse de 60 points supplémentaires jusqu'à l'utilisation totale, et le
taux servi aux fournisseurs qui en découle avec un facteur de réserve de 10~\% (\cref{prop:m3:defi-credit-and-leverage:supply}). Paramètres illustratifs de la
forme à deux pentes. Données~: le tutoriel du chapitre.}\label{fig:m3:defi-credit-and-leverage:rates}
\end{omfigure}

\begin{dated}{2026-09}{Modèle de taux d'un protocole de prêt et frais de flash loan}\label{dat:m3:defi-credit-and-leverage:aave}
La documentation d'Aave décrit sa stratégie de taux de la v3 comme reposant sur deux pentes, l'une sous un ratio d'utilisation optimal et l'autre à partir
de ce point jusqu'à 100~\%, avec un taux d'emprunt variable de base, par réserve. Les positions sont surcollatéralisées, le risque est suivi au moyen d'un \omterm{def:m3:defi-credit-and-leverage:ltv}{facteur de santé}
avec des seuils de liquidation par réserve, et une position dont le \omterm{def:m3:defi-credit-and-leverage:ltv}{facteur de santé} est inférieur à 1 peut être liquidée par n'importe qui, qui rembourse une partie de la
dette et reçoit du collatéral avec une remise, le \omterm{def:m3:defi-credit-and-leverage:ltv}{bonus de liquidation}. Dans les contrats v3, une liquidation peut rembourser au plus 50~\% de la dette
(le facteur de clôture), et la totalité lorsque le \omterm{def:m3:defi-credit-and-leverage:ltv}{facteur de santé} est égal ou inférieur à 0.95 ou que la position est inférieure à 2\,000~dollars. Les frais de \omterm{def:m3:defi-credit-and-leverage:flash}{flash loan} ont été fixés à 0.05~\%
lors du déploiement et peuvent être modifiés par la gouvernance.
\end{dated}

\section{Liquidations}

\begin{definition}[Ratio prêt/valeur, facteur de santé, bonus de liquidation]\label{def:m3:defi-credit-and-leverage:ltv}
Le \emph{ratio prêt/valeur}\index{ratio prêt/valeur} d'un actif de collatéral est le montant maximal qu'un emprunteur peut emprunter par dollar de cet actif. Le
\emph{facteur de santé}\index{facteur de santé} d'un compte est la somme, sur son collatéral, de la valeur multipliée par le seuil de liquidation, divisée par la
valeur de sa dette~; inférieur à un, le compte peut être liquidé. Le \emph{bonus de liquidation}\index{bonus de liquidation} est la remise à laquelle un
liquidateur reçoit du collatéral en échange de la dette qu'il rembourse.
\end{definition}

Le \omterm{def:m3:defi-credit-and-leverage:ltv}{ratio prêt/valeur} est une décote (livre~1, chapitre~6) appliquée à l'emprunt~; le seuil de liquidation, un peu plus élevé, est la
décote à laquelle la position est reprise. Il n'y a pas d'appel de marge~: l'emprunteur n'est pas sommé de fournir davantage de collatéral, il est liquidé par celui qui
agit le premier, et le bonus le rémunère pour agir. Un emprunteur disposant de 100 ether de collatéral, avec un seuil de 82.5~\% et 240\,000~dollars de dette, a un
\omterm{def:m3:defi-credit-and-leverage:ltv}{facteur de santé} de 1.03 pour un ether à 3\,000~dollars et est liquidé en dessous de 2\,909.

\begin{proposition}[Arithmétique de la liquidation]\label{prop:m3:defi-credit-and-leverage:liq}
Un liquidateur rembourse un montant $R$ de dette, au plus le facteur de clôture $c$ multiplié par la dette, et reçoit du collatéral valant $R(1 + \beta)$ au
prix de l'\omterm{def:m3:blockchains-for-traders:contract}{oracle}, avec $\beta$ le bonus. Le \omterm{def:m3:defi-credit-and-leverage:ltv}{facteur de santé} de l'emprunteur après une liquidation de $R$ est
$(C_\theta - R(1 + \beta)\theta)/(D - R)$, où $C_\theta$ est la valeur du collatéral pondérée par les seuils, $D$ la valeur de la dette et $\theta$ le
seuil du collatéral~; il ne croît que si $(1 + \beta)\theta < C_\theta/D$, ce qui est vérifié lorsque le \omterm{def:m3:defi-credit-and-leverage:ltv}{facteur de santé} était supérieur à $(1 + \beta)\theta$.
\end{proposition}

\begin{proof}
La dette diminue de $R$ et le collatéral pondéré par les seuils diminue de $R(1 + \beta)\theta$. Le rapport $(C_\theta - xR)/(D - R)$ avec $x = (1 +
\beta)\theta$ dépasse $C_\theta/D$ exactement lorsque $C_\theta/D > x$.
\end{proof}

\begin{omfigure}
\begin{tikzpicture}
\begin{axis}[width=11cm,height=5cm,x dir=reverse,
  xlabel={prix de l'ether, en dollars},ylabel={facteur de santé},
  tick label style={font=\small},label style={font=\small},
  xmin=2000,xmax=3300,ymin=0.6,ymax=1.2,grid=major,grid style={gray!25},
  xtick={2000,2250,2500,2750,3000,3250},xticklabels={2\,000,2\,250,2\,500,2\,750,3\,000,3\,250},scaled x ticks=false,
  table/col sep=comma]
\addplot[omThm,very thick] table[x=price,y=hf]{figdata/markets-3/23-defi-credit-and-leverage/health.csv};
\draw[omDef,thick,dashed] (axis cs:3300,1) -- (axis cs:2000,1);
\draw[omExo,thick] (axis cs:2909,0.6) -- (axis cs:2909,1.2);
\node[font=\footnotesize,anchor=west,text=omExo] at (axis cs:2890,0.68) {liquidation en dessous de 2\,909};
\end{axis}
\end{tikzpicture}

\omcaption{\omterm{def:m3:defi-credit-and-leverage:ltv}{Facteur de santé} d'un emprunteur disposant de 100 ether de collatéral (seuil de 82.5~\%) et de 240\,000~dollars de dette, à mesure que le prix de l'ether baisse.
Données~: le tutoriel du chapitre.}\label{fig:m3:defi-credit-and-leverage:health}
\end{omfigure}

La proposition réserve une mauvaise surprise. Une position qui tombe bien en dessous de un avant que quiconque n'agisse, parce que le cours a fait un saut ou que la chaîne était saturée, peut
être enfoncée davantage par sa propre liquidation~: chaque dollar remboursé retire plus d'un dollar de collatéral pondéré par les seuils. Lorsque le collatéral restant
vaut moins que la dette, le pool se retrouve avec une dette irrécouvrable que ses fournisseurs supportent.

\section{Prêts flash}

\begin{definition}[Flash loan (prêt flash)]\label{def:m3:defi-credit-and-leverage:flash}
Un \emph{flash loan}\index{flash loan} est un prêt sans collatéral qui doit être remboursé, avec des frais, dans la même transaction~; si ce n'est pas
le cas, l'ensemble de la transaction est annulé et le prêt n'a jamais eu lieu.
\end{definition}

Un \omterm{def:m3:defi-credit-and-leverage:flash}{flash loan} ne présente aucun risque de crédit pour le prêteur, car l'atomicité de la chaîne impose le remboursement~: la transaction se termine soit avec le pool remboursé
ou elle n'existe pas. Il permet à quiconque d'agir avec le capital d'un pool le temps d'une seule transaction. L'usage canonique est une liquidation par un liquidateur sans capital~:
emprunter les \omterm{def:m3:blockchains-for-traders:stable}{stablecoins} pour rembourser, recevoir le collatéral avec son bonus, le vendre sur un \omterm{def:m3:automated-market-makers:amm}{teneur de marché automatisé}, rembourser le \omterm{def:m3:defi-credit-and-leverage:flash}{flash loan} et garder le reste.
Dans le tutoriel, la liquidation de la moitié de la dette de l'emprunteur à un cours de l'ether de 2\,800 saisit 45 ether, les vend pour 124\,505~USD dans un pool de
5\,000 ether, rembourse les 120\,000~USD empruntés et 60~USD de frais, et laisse 4\,445~USD. La même atomicité sert aux attaques~: un \omterm{def:m3:defi-credit-and-leverage:flash}{flash loan} peut faire bouger un cours,
utiliser le cours ainsi déplacé et le rétablir, le tout dans une seule transaction. Qin, Zhou, Livshits et Gervais ont analysé deux de ces attaques, toutes deux
en février 2020, et ont mesuré des rendements sur investissement supérieurs à 500\,000~\%.

\begin{omfigure}
\centering
\begin{tikzpicture}[font=\small,
  s/.style={draw,thick,rounded corners=2pt,align=center,minimum height=0.9cm,text width=2.2cm},
  arr/.style={-{Stealth[length=2.2mm]},thick}]
\node[s,draw=omDef,fill=omDef!8] (a) at (0,0) {emprunt flash\\120\,000 USDC};
\node[s,draw=omThm,fill=omThm!8] (b) at (2.8,0) {rembourser la moitié\\de la dette};
\node[s,draw=omThm,fill=omThm!8] (c) at (5.6,0) {recevoir 45 ETH\\(bonus de 5~\%)};
\node[s,draw=omProp,fill=omProp!8] (d) at (8.4,0) {vendre dans un pool~:\\124\,505 USDC};
\node[s,draw=omDef,fill=omDef!8] (e) at (11.2,0) {rembourser 120\,060~;\\garder 4\,445};
\draw[arr] (a) -- (b); \draw[arr] (b) -- (c); \draw[arr] (c) -- (d); \draw[arr] (d) -- (e);
\draw[thick,gray,dashed,rounded corners] (-1.3,-0.8) rectangle (12.5,0.8);
\node[font=\footnotesize,gray,anchor=north] at (5.6,-0.85) {une seule transaction~: si la dernière étape échoue, rien de tout cela n'a eu lieu};
\end{tikzpicture}

\omcaption{Une liquidation financée par un \omterm{def:m3:defi-credit-and-leverage:flash}{flash loan}, avec les chiffres du tutoriel. Schématique.}\label{fig:m3:defi-credit-and-leverage:flash}
\end{omfigure}

\section{Staking liquide, restaking et rendement de staking}

\begin{definition}[Rendement de staking, jeton de staking liquide, restaking]\label{def:m3:defi-credit-and-leverage:staking}
Le \emph{rendement de staking}\index{rendement de staking} est le rendement obtenu en immobilisant le \omterm{def:m3:blockchains-for-traders:key}{jeton} d'une chaîne en \omterm{def:m3:blockchains-for-traders:pow}{preuve d'enjeu} auprès d'un
\omterm{def:m3:blockchains-for-traders:pow}{validateur} (\cref{ch:m3:blockchains-for-traders})~: émissions nouvelles et une part des frais, diminuées des pénalités. Un \emph{jeton de staking liquide}\index{jeton de staking liquide} est un
\omterm{def:m3:blockchains-for-traders:key}{jeton} transférable qui représente des \omterm{def:m3:blockchains-for-traders:key}{jetons} stakés et leurs récompenses courues, émis par un protocole qui mise les dépôts pour le compte de ses utilisateurs.
Le \emph{restaking}\index{restaking} consiste à engager des \omterm{def:m3:blockchains-for-traders:key}{jetons} déjà stakés, ou des \omterm{def:m3:blockchains-for-traders:key}{jetons} de staking liquide, en garantie de services supplémentaires, pour des
récompenses additionnelles et au prix d'un risque supplémentaire de les voir détruits en cas de comportement fautif.
\end{definition}

\begin{dated}{2026-09}{Un rendement de staking}\label{dat:m3:defi-credit-and-leverage:lido}
L'API publique de Lido a fait état d'un taux annuel en pourcentage moyen mobile sur sept jours de 2.24~\% pour son \omterm{def:m3:blockchains-for-traders:key}{jeton} d'ether staké, stETH, le 24 septembre 2026,
et 9.76 millions d'ethers stakés par son intermédiaire, soit environ 22~\% des 43.8 millions d'ethers détenus par les \omterm{def:m3:blockchains-for-traders:pow}{validateurs} sur la beacon chain ce jour-là.
\end{dated}

Le \omterm{def:m3:defi-credit-and-leverage:staking}{rendement de staking} est ce qui se rapproche le plus, dans la cryptographie, d'un taux sans risque pour l'ether, et le \omterm{def:m3:defi-credit-and-leverage:staking}{jeton de staking liquide} constitue le collatéral sur lequel
repose une grande part de l'effet de levier~: on dépose du stETH, on emprunte de l'ether, on le stake, et on recommence, en touchant le \omterm{def:m3:defi-credit-and-leverage:staking}{rendement de staking} sur plusieurs fois le capital tant que le taux d'emprunt
reste inférieur à celui-ci. Les risques de cette boucle sont ceux de tout carry trade financé à court terme~: le taux d'emprunt peut dépasser le rendement lorsque le taux d'utilisation augmente,
et le \omterm{def:m3:defi-credit-and-leverage:staking}{jeton de staking liquide} peut se négocier en dessous de la valeur de l'ether qu'il représente lorsque les détenteurs veulent sortir plus vite que ne le permettent les retraits du staking,
ce qui déclenche des liquidations chez les boucleurs. Le \omterm{def:m3:defi-credit-and-leverage:staking}{restaking} superpose d'autres créances sur le même collatéral~: les contrats d'EigenLayer acceptent des \omterm{def:m3:defi-credit-and-leverage:staking}{jetons de staking liquide},
de l'ether de la beacon chain et d'autres \omterm{def:m3:blockchains-for-traders:key}{jetons}, et laissent chaque opérateur allouer une part de son enjeu délégué à un service donné, qui peut le saisir en cas de faute.

\section{Stablecoins algorithmiques et manipulation d'oracle}

\begin{definition}[Stablecoin algorithmique]\label{def:m3:defi-credit-and-leverage:algo}
Un \emph{stablecoin algorithmique}\index{stablecoin algorithmique} est un \omterm{def:m3:blockchains-for-traders:key}{jeton} qui vise une valeur fixe sans détenir de réserves de cette valeur, en s'appuyant
plutôt sur un mécanisme qui l'échange contre un autre \omterm{def:m3:blockchains-for-traders:key}{jeton} à valeur flottante, créé ou détruit pour absorber la demande.
\end{definition}

Le mécanisme ne tient l'ancrage que tant que le \omterm{def:m3:blockchains-for-traders:key}{jeton} flottant vaut plus que les \omterm{def:m3:blockchains-for-traders:stable}{stablecoins} en circulation~; lorsque la confiance s'effondre, les rachats créent
des \omterm{def:m3:blockchains-for-traders:key}{jetons} flottants plus vite que quiconque n'en veut, leur cours baisse, et chaque rachat en crée davantage. Selon la SEC, Terraform Labs
a commercialisé son \omterm{def:m3:blockchains-for-traders:stable}{stablecoin} UST, qui devait conserver son ancrage au dollar en étant échangeable contre son \omterm{def:m3:blockchains-for-traders:key}{jeton} LUNA, comme «~à rendement~», versant jusqu'à 20~\%
par l'intermédiaire de son protocole Anchor~; en mai 2022, UST a perdu son ancrage, et lui et ses \omterm{def:m3:blockchains-for-traders:key}{jetons} sœurs sont tombés à près de zéro.

\begin{definition}[Manipulation d'oracle]\label{def:m3:defi-credit-and-leverage:oracle}
La \emph{manipulation d'oracle}\index{manipulation d'oracle} consiste à faire bouger le cours qu'un \omterm{def:m3:blockchains-for-traders:contract}{oracle} transmet à un \omterm{def:m3:blockchains-for-traders:contract}{smart contract}, généralement en négociant
sur les marchés peu profonds qu'il lit, afin d'emprunter contre un collatéral gonflé, d'éviter une liquidation ou d'en déclencher une.
\end{definition}

\begin{dated}{2026-09}{L'affaire Mango Markets}\label{dat:m3:defi-credit-and-leverage:mango}
La CFTC a allégué en janvier 2023 que, le 11 octobre 2022, Avraham Eisenberg, par l'intermédiaire de deux comptes sur Mango Markets, a constitué de grosses positions à effet de levier
sur un swap portant sur le \omterm{def:m3:blockchains-for-traders:key}{jeton} MNGO, a fait grimper le cours de MNGO sur trois bourses alimentant l'\omterm{def:m3:blockchains-for-traders:contract}{oracle} de la plateforme, lequel a fait état d'une hausse de plus de treize fois en 30 minutes,
et a utilisé les positions gonflées comme collatéral pour retirer plus de 110~millions de dollars~USD, dont environ 67~millions de dollars~USD ont été restitués par la suite. Un jury l'a déclaré coupable de
fraude sur les matières premières, de manipulation sur les matières premières et de fraude électronique~; le 23 mai 2025, le juge du procès a accueilli la requête de l'intéressé fondée sur la Rule 29, annulant les deux premiers
chefs d'accusation et prononçant l'acquittement sur le troisième.
\end{dated}

\begin{proposition}[La hausse qui vide un pool]\label{prop:m3:defi-credit-and-leverage:pump}
Un pool prête jusqu'à $L$ contre un collatéral de $q$ \omterm{def:m3:blockchains-for-traders:key}{jetons} au cours $P$ et au \omterm{def:m3:defi-credit-and-leverage:ltv}{ratio prêt/valeur} $\lambda$, le cours étant lu sur un marché à produit constant
qui détient $y$ dollars. Faire monter le cours de ce marché d'un facteur $k$ coûte $y(\sqrt{k} - 1)$ dollars et achète une fraction $1 - 1/\sqrt{k}$ de ses \omterm{def:m3:blockchains-for-traders:key}{jetons}.
L'attaquant peut emprunter la totalité du pool dès que $k \ge k^* = L/(\lambda qP)$~; son gain, en empruntant puis en abandonnant le collatéral, est
$\min(L, \lambda qPk) - qP - y(\sqrt{k} - 1)$, auquel s'ajoute la valeur des \omterm{def:m3:blockchains-for-traders:key}{jetons} achetés ensuite.
\end{proposition}

\begin{proof}
En maintenant $xy$ constant pendant que le cours $y/x$ monte d'un facteur $k$, $y$ est multiplié par $\sqrt k$ et $x$ divisé par $\sqrt k$. Le pouvoir d'emprunt vaut $\lambda qPk$,
plafonné par le pool~; le collatéral reste en place.
\end{proof}

\begin{omfigure}
\begin{tikzpicture}
\begin{axis}[width=11cm,height=5cm,
  xlabel={hausse du cours source de l'oracle, multiple},ylabel={gain de l'attaquant, millions de USD},
  tick label style={font=\small},label style={font=\small},
  xmin=1,xmax=40,ymin=-5,ymax=45,grid=major,grid style={gray!25},
  table/col sep=comma]
\addplot[omThm,very thick] table[x=k,y=gain]{figdata/markets-3/23-defi-credit-and-leverage/pump.csv};
\draw[omDef,thick,dashed] (axis cs:16.67,-5) -- (axis cs:16.67,45);
\node[font=\footnotesize,anchor=west,text=omDef] at (axis cs:17.2,5) {$k^* = 16.7$~: pool vidé par l'emprunt};
\end{axis}
\end{tikzpicture}

\omcaption{La hausse de l'\omterm{def:m3:blockchains-for-traders:contract}{oracle} de la \cref{prop:m3:defi-credit-and-leverage:pump}~: 5 millions de \omterm{def:m3:blockchains-for-traders:key}{jetons} à 1~USD en collatéral (\omterm{def:m3:defi-credit-and-leverage:ltv}{ratio prêt/valeur} de 60~\%), un pool de 50
millions de dollars~USD, et un marché source détenant 2~millions de dollars~USD~; les \omterm{def:m3:blockchains-for-traders:key}{jetons} achetés pendant la hausse sont supposés sans valeur par la suite. Le gain devient positif pour une hausse
inférieure à deux fois, atteint son maximum lorsque le pool est vide, puis décroît, car toute hausse supplémentaire ne fait qu'ajouter des coûts. Paramètres illustratifs. Données~: le tutoriel du chapitre.}\label{fig:m3:defi-credit-and-leverage:pump}
\end{omfigure}

Les défenses découlent de la formule~: des \omterm{def:m3:defi-credit-and-leverage:ltv}{ratios prêt/valeur} qui baissent à mesure que les marchés du collatéral deviennent moins profonds, des caps sur le montant qui peut être emprunté
pour un même collatéral, des \omterm{def:m3:blockchains-for-traders:contract}{oracles} qui lisent des marchés profonds et lissent dans le temps, et des coupe-circuits face aux brusques variations de prix. Aucun n'est gratuit~:
chacune rend le \omterm{def:m3:defi-credit-and-leverage:pool}{pool de prêt} moins utile aux emprunteurs honnêtes.

\section{Tutoriel~: un pool de prêt en période de tension}\label{tut:m3:defi-credit-and-leverage:pool}

\begin{tutorial}
\textbf{Objectif.} Simuler un \omterm{def:m3:defi-credit-and-leverage:pool}{pool de prêt}~: le taux d'emprunt sous forme de courbe coudée en fonction du taux d'utilisation, le \omterm{def:m3:defi-credit-and-leverage:ltv}{facteur de santé} d'un emprunteur lors d'une baisse de prix, une
liquidation avec facteur de clôture et bonus, et une liquidation financée par un \omterm{def:m3:defi-credit-and-leverage:flash}{flash loan}. \textbf{État final~:} les figures du chapitre et les chiffres du
texte.
\begin{tutsteps}
\item \textbf{Liquidation et \omterm{def:m3:defi-credit-and-leverage:flash}{flash loan}.} Les deux fonctions les plus délicates du moteur.
\omcode{code/firm/lendpool/firm_lendpool.py}{93}{119}{Liquidation avec facteur de clôture et bonus, et un flash loan qui est annulé s'il n'est pas remboursé.}\label{lst:m3:defi-credit-and-leverage:liq}
\item \textbf{La liquidation par \omterm{def:m3:defi-credit-and-leverage:flash}{flash loan}.} Emprunter, liquider, vendre le collatéral dans un pool avec \texttt{firm.amm}, rembourser.
\omcode{code/markets-3/23-defi-credit-and-leverage/python/m3_defi.py}{38}{57}{Une liquidation financée par un flash loan.}\label{lst:m3:defi-credit-and-leverage:flash}
\item \textbf{Exécuter} \texttt{flash\_liquidation()}, \texttt{pump\_attack} sur des hausses de 1 à 40, et \texttt{fig\_defi.py}.
\end{tutsteps}
\textbf{Pour aller plus loin.} Faire chuter brutalement le prix à 2\,400 avant la liquidation et déterminer la mauvaise créance du pool~; ajouter un second liquidateur en concurrence dans le même
bloc~; faire de l'\omterm{def:m3:blockchains-for-traders:contract}{oracle} une moyenne sur 30~minutes et recalculer le coût de la hausse.
\end{tutorial}

\section{Construction~: le moteur du pool de prêt}\label{bld:m3:defi-credit-and-leverage:lendpool}

\begin{build}
\bfield{Rôle} La firme miniature prête, emprunte et liquide on-chain~: elle doit connaître ses \omterm{def:m3:defi-credit-and-leverage:ltv}{facteurs de santé}, les taux qu'elle paiera à mesure que le taux d'utilisation évolue,
le profit d'une liquidation après la vente du collatéral, et l'exposition d'un pool qu'elle alimente à la mauvaise créance.
\bfield{Interface} \texttt{RateModel(base, slope1, slope2, optimal)}~; \texttt{Reserve} avec \texttt{utilisation}, \texttt{rates}, \texttt{accrue},
\texttt{available}~; \texttt{Pool} avec \texttt{supply}, \texttt{borrow}, \texttt{health}, \texttt{liquidate}, \texttt{flash\_loan}. Ventes de collatéral via
\texttt{firm.amm}.
\bfield{Règles} Dettes et dépôts tenus en soldes normalisés par des indices~; emprunt refusé au-delà du \omterm{def:m3:defi-credit-and-leverage:ltv}{ratio prêt/valeur} ou de la liquidité disponible~; liquidation
seulement si le \omterm{def:m3:defi-credit-and-leverage:ltv}{facteur de santé} est inférieur à un et dans la limite du facteur de clôture~; un \omterm{def:m3:defi-credit-and-leverage:flash}{flash loan} non remboursé lève une exception et ne laisse aucune trace.
\bfield{Tests d'acceptation} \texttt{code/firm/lendpool/tests/}~: le coude de la courbe~; \omterm{def:m3:defi-credit-and-leverage:ltv}{facteur de santé} et accumulation des intérêts~; refus au-delà du \omterm{def:m3:defi-credit-and-leverage:ltv}{ratio prêt/valeur}~; refus de liquider un
compte sain~; saisie avec le bonus~; frais d'un \omterm{def:m3:defi-credit-and-leverage:flash}{flash loan} et son annulation.
\bfield{Pour aller plus loin} Plusieurs actifs en collatéral avec des seuils différents~; comptabilité des mauvaises créances~; modes d'isolement et caps d'emprunt~; la réaction du taux à une ruée sur
les dépôts.
\end{build}

\begin{omsources}
\item Documentation d'Aave~: v3 overview, interest-rate strategy, flash loans~; le contrat LiquidationLogic du dépôt aave-v3-origin. Consulté en septembre 2026.
\item EigenLayer core contracts documentation (Layr-Labs/eigenlayer-contracts, docs/README), septembre 2026.
\item K.~Qin, L.~Zhou, B.~Livshits et A.~Gervais, ``Attacking the DeFi Ecosystem with Flash Loans for Fun and Profit'', Financial Cryptography 2021
(arXiv:2003.03810).
\item CFTC, communiqué de presse 8647-23 (charges against Avraham Eisenberg), janvier 2023~; United States v.\ Eisenberg (S.D.N.Y.), avis et ordonnance du 23
mai 2025 (Doc 220).
\item SEC, communiqué de presse 2023-32, ``SEC Charges Terraform and CEO Do Kwon with Defrauding Investors'', 16 février 2023.
\item Lido, public stETH APR and statistics endpoints~; beacon-chain validator balances from a public beacon node~; 24 septembre 2026.
\end{omsources}

\section{Exercices}

\begin{exercise}[$\star$]\label{exo:m3:defi-credit-and-leverage:1}
Avec la courbe de taux du chapitre, déterminer les taux d'emprunt et de dépôt pour des taux d'utilisation de 50~\% et de 95~\%.
\end{exercise}

\begin{exercise}[$\star$]\label{exo:m3:defi-credit-and-leverage:2}
Un emprunteur détient 50~ether en collatéral (seuil de 82.5~\%) et une dette de 100\,000~USD. Déterminer le prix de l'ether auquel il est liquidé.
\end{exercise}

\begin{exercise}[$\star$]\label{exo:m3:defi-credit-and-leverage:3}
Expliquer pourquoi un \omterm{def:m3:defi-credit-and-leverage:flash}{flash loan} ne présente aucun risque de crédit pour le pool.
\end{exercise}

\begin{exercise}[$\star\star$]\label{exo:m3:defi-credit-and-leverage:4}
Un liquidateur rembourse 50\,000~USD de dette contre de l'ether à 2\,500, avec un bonus de 5~\%. Déterminer la quantité d'ether qu'il reçoit.
\end{exercise}

\begin{exercise}[$\star\star$]\label{exo:m3:defi-credit-and-leverage:5}
Un emprunteur a un \omterm{def:m3:defi-credit-and-leverage:ltv}{facteur de santé} de 0.80, le seuil du collatéral est de 82.5~\% et le bonus de 10~\%. Déterminer si la liquidation améliore son \omterm{def:m3:defi-credit-and-leverage:ltv}{facteur de santé}.
\end{exercise}

\begin{exercise}[$\star\star$]\label{exo:m3:defi-credit-and-leverage:6}
Expliquer la stratégie de looping sur un \omterm{def:m3:defi-credit-and-leverage:staking}{jeton de staking liquide} et citer deux façons dont elle fait perdre de l'argent.
\end{exercise}

\begin{exercise}[$\star\star\star$]\label{exo:m3:defi-credit-and-leverage:7}
\emph{Programmation.} À l'aide de \texttt{pump\_attack}, déterminer la plus petite hausse à partir de laquelle l'attaque est rentable, avec et sans que les \omterm{def:m3:blockchains-for-traders:key}{jetons} achetés conservent une valeur de 1~USD.
\end{exercise}

\begin{exercise}[$\star\star\star$]\label{exo:m3:defi-credit-and-leverage:8}
\emph{Trouver la faille.} «~Notre pool est sûr~: chaque prêt est surcollatéralisé.~»
\end{exercise}

\section{Problème~: la hausse de l'oracle}

\begin{problem}[{Problème du week-end --- vider un pool par l'emprunt}]\label{pb:m3:defi-credit-and-leverage:1}
Un \omterm{def:m3:defi-credit-and-leverage:pool}{pool de prêt} détient 50~millions de dollars de \omterm{def:m3:blockchains-for-traders:stable}{stablecoins} et accepte un \omterm{def:m3:blockchains-for-traders:key}{jeton} peu liquide comme collatéral, avec un \omterm{def:m3:defi-credit-and-leverage:ltv}{ratio prêt/valeur} de 60~\%, évalué par un \omterm{def:m3:blockchains-for-traders:contract}{oracle} qui lit un
marché à produit constant détenant 2~millions de dollars et 2~millions de \omterm{def:m3:blockchains-for-traders:key}{jetons} (prix de 1~USD). Un attaquant détient 5~millions de \omterm{def:m3:blockchains-for-traders:key}{jetons}.

\medskip\noindent\textbf{Partie I --- Le mécanisme.}
\begin{enumerate}
\item Quel montant l'attaquant peut-il emprunter sans manipulation~?
\item De quel facteur le prix de l'\omterm{def:m3:blockchains-for-traders:contract}{oracle} doit-il augmenter pour que l'attaquant puisse emprunter tout le pool~?
\item Combien coûte le fait de faire monter le prix du marché de ce facteur, et combien de \omterm{def:m3:blockchains-for-traders:key}{jetons} cela permet-il d'acheter~?
\item Quel est le gain de l'attaquant si les \omterm{def:m3:blockchains-for-traders:key}{jetons} achetés ne valent plus rien~? Et s'ils conservent 1~USD~?
\item Pourquoi le gain diminue-t-il pour les hausses supérieures à ce facteur~?
\end{enumerate}

\medskip\noindent\textbf{Partie II --- Les seuils.}
\begin{enumerate}[resume]
\item Quelle est la plus petite hausse rentable~?
\item Comment la réponse change-t-elle si le marché source détient 20~millions de dollars~?
\item Et si le \omterm{def:m3:defi-credit-and-leverage:ltv}{ratio prêt/valeur} est de 30~\%~?
\item Quel cap d'emprunt sur ce collatéral limiterait à zéro le gain de l'attaquant, quelle que soit la hausse~?
\item Que changerait un \omterm{def:m3:blockchains-for-traders:contract}{oracle} moyenné dans le temps sur 30~minutes~?
\end{enumerate}

\medskip\noindent\textbf{Partie III --- L'affaire.}
\begin{enumerate}[resume]
\item Qu'a allégué la CFTC au sujet de ce qui s'est produit sur Mango Markets~?
\item Qu'a décidé le juge de première instance après le verdict du jury, et sur quel type de requête~?
\item Qui supporte la perte lorsqu'un pool est vidé par l'emprunt~?
\item En quoi cette attaque diffère-t-elle d'une attaque par \omterm{def:m3:defi-credit-and-leverage:flash}{flash loan}~?
\item Que ferait ensuite la gouvernance d'un protocole de prêt~?
\end{enumerate}

\medskip\noindent\textbf{Partie IV --- Jugement.}
\begin{enumerate}[resume]
\item Le dysfonctionnement vient-il de la conception du pool ou du comportement de l'attaquant~?
\item Un pool doit-il accepter des \omterm{def:m3:blockchains-for-traders:key}{jetons} peu liquides comme collatéral~?
\item Comment un fournisseur doit-il choisir les pools auxquels il fournit des liquidités~?
\item Énoncer le \emph{résultat nommé}~: la manipulation qui vide le pool de ses actifs, et le coût de l'attaquant comparé au gain.
\item En une phrase~: qu'est-ce qu'un \omterm{def:m3:blockchains-for-traders:contract}{oracle} pour un \omterm{def:m3:defi-credit-and-leverage:pool}{pool de prêt}~?
\end{enumerate}
\end{problem}

\section{Questions d'entretien}

\begin{interviewq}[$\star$ \iqroles{trader}]\label{iq:m3:defi-credit-and-leverage:1}
Qu'est-ce qu'un \omterm{def:m3:defi-credit-and-leverage:ltv}{facteur de santé}, et que se passe-t-il lorsqu'il passe sous un~?
\end{interviewq}

\begin{interviewq}[$\star$ \iqroles{développeur}]\label{iq:m3:defi-credit-and-leverage:2}
Comment fonctionne un \omterm{def:m3:defi-credit-and-leverage:flash}{flash loan}, et pourquoi ne peut-il exister que sur une \omterm{def:m3:blockchains-for-traders:chain}{blockchain}~?
\end{interviewq}

\begin{interviewq}[$\star\star$ \iqroles{chercheur}]\label{iq:m3:defi-credit-and-leverage:3}
Pourquoi les \omterm{def:m3:defi-credit-and-leverage:pool}{pools de prêt} utilisent-ils une \omterm{def:m3:defi-credit-and-leverage:pool}{courbe de taux coudée} plutôt qu'une courbe linéaire~?
\end{interviewq}

\begin{interviewq}[$\star\star$ \iqroles{risque}]\label{iq:m3:defi-credit-and-leverage:4}
Comment fixeriez-vous le \omterm{def:m3:defi-credit-and-leverage:ltv}{ratio prêt/valeur} d'un nouveau \omterm{def:m3:blockchains-for-traders:key}{jeton} de collatéral~?
\end{interviewq}

\begin{interviewq}[$\star\star$ \iqroles{trader, chercheur}]\label{iq:m3:defi-credit-and-leverage:5}
Pourquoi un \omterm{def:m3:defi-credit-and-leverage:algo}{stablecoin algorithmique} versant 20~\% s'est-il effondré, et qu'auriez-vous surveillé~?
\end{interviewq}

\begin{interviewq}[$\star\star\star$ \iqroles{développeur, risque}]\label{iq:m3:defi-credit-and-leverage:6}
Concevez un bot de liquidation qui se fait concurrence de manière rentable sans prendre de risque de stock défavorable.
\end{interviewq}
